\section{Source Subsystem Detail and Tables}\label{app:source-03}
\noindent\textit{Historical source detail. Several numerical assumptions and design descriptions in this transcription have been superseded by Sections 1--5 and Appendix L. These tables are retained for traceability and are not the current approved mission baseline.}

\noindent\textit{The following is retained from the original project draft for supporting detail and comparison. Numerical claims are preliminary unless specifically reconciled in the main text.}

The ANTARA spacecraft architecture has been developed as a modular,
scalable, and mission-oriented CubeSat platform optimised for scientific
radiation measurements withi n the South Atlantic Anomaly. The
spacecraft follows a centralised avionics architecture in which the
On-Board Computer (OBC) functions as the primary command. The
communications subsystem, Payload Subsystem and ADCS subsystem work on
the basis of the instantaneous mode. Following is the spacecraft system
architecture.

\subsection{System Architecture Overview}

The ANTARA spacecraft is a \textbf{3U CubeSat}, conforming to the
standard 3U form factor of 100 × 100 × 340.5 mm as defined by the
CubeSat Design Specification (CDS) within the standard 3U mass
allowance. Compliance with the applicable CDS revision, including
deployer-imposed mass, center-of-gravity, and rail/standoff constraints,
will be formally verified once a launch integrator and deployer are
selected.

The five primary spacecraft subsystems integrated around a centralised
avionics architecture are: On-Board Computer (OBC), Electrical Power
System (EPS), Attitude Determination and Control System (ADCS),
Communications System (COMMS), and Scientific Payload.

\subsubsection{Electrical Power System (EPS)}

The Electrical Power Subsystem is responsible for generating, storing,
and distributing electrical power to all spacecraft loads. The COTS
model we have taken for the simulation is the ISISPACE ICEPS2 Type A
compact EPS, which includes an integrated 2-cell battery pack,
autonomous battery heating, I2C communication, up to 6 output voltage
domains, and protection features including overcurrent and
reverse-current protection.

\begin{longtable}{@{}
  >{\raggedright\arraybackslash}p{(\columnwidth - 2\tabcolsep) * \real{0.3333}}
  >{\raggedright\arraybackslash}p{(\columnwidth - 2\tabcolsep) * \real{0.6667}}@{}}
\caption{ICEPS2 Type A parameters}\\
\toprule\noalign{}
\textbf{Parameter} & \textbf{Preliminary Requirement} \\
Mass & 184 ± 5 g \\
Volume & 96 × 92 × 26.45 mm \\
Energy storage & 22.5 Wh (Type A specification) \\
Operating temperature & −20°C to +70°C \\
Interfaces & I2C housekeeping/control \\
Voltage domains & Regulated 3.3 V, 5 V, unregulated bus; up to 6 output domains \\
\midrule\noalign{}
\endhead
\bottomrule\noalign{}
\endlastfoot
\end{longtable}

Note: The mission simulation uses a 25 Wh battery model as a
conservative planning estimate. The final power budget must be
reconciled against the ICEPS2 Type A 22.5 Wh specification at PDR.

\textbf{Mode Power Estimates:}

\begin{longtable}{@{}
  >{\raggedright\arraybackslash}p{(\columnwidth - 4\tabcolsep) * \real{0.2543}}
  >{\raggedright\arraybackslash}p{(\columnwidth - 4\tabcolsep) * \real{0.4837}}
  >{\raggedright\arraybackslash}p{(\columnwidth - 4\tabcolsep) * \real{0.2620}}@{}}
\caption{Spacecraft operating-mode power estimates}\\
\toprule\noalign{}
\textbf{Mode} & \textbf{Active Components} & \textbf{Estimated Power (W)} \\
Cold \& Dark Boot & EPS + OBC & 0.8 \\
Detumbling & EPS + OBC + ADCS MTQs & 3--6 \\
Safe Mode & EPS + OBC + Radio RX & 1--2 \\
Sun-Pointing & EPS + OBC + ADCS & 2--3 \\
Detection Mode & EPS + OBC + Payload + ADCS & 6--9 \\
Relay Mode & EPS + OBC + Radio TX + ADCS & 4--7 \\
\midrule\noalign{}
\endhead
\bottomrule\noalign{}
\endlastfoot
\end{longtable}

\begin{longtable}{@{}
  >{\raggedright\arraybackslash}p{(\columnwidth - 6\tabcolsep) * \real{0.1940}}
  >{\raggedright\arraybackslash}p{(\columnwidth - 6\tabcolsep) * \real{0.2587}}
  >{\raggedright\arraybackslash}p{(\columnwidth - 6\tabcolsep) * \real{0.3035}}
  >{\raggedright\arraybackslash}p{(\columnwidth - 6\tabcolsep) * \real{0.2438}}@{}}
\caption{Preliminary subsystem electrical power budget}\\
\toprule\noalign{}
\textbf{Subsystem} & \textbf{Idle / Standby (W)} & \textbf{Nominal (W)} & \textbf{Peak (W)} \\
EPS & -0.093 & 0 & 0 \\
Battery & - & charging capacity(28.8Wh) & - \\
OBC & - & -0.17 & -0.9 \\
Radio & -0.15 (rx) & -2.64 (tx) & -3 (TX max) \\
Unidirectional antenna & -0.04 & -0.06 & -2 (Deployment) \\
ADCS & -0.8 & -2.6 & -4.5 \\
Payload

(Assumed) & - & -3.0 & -5.0 \\
Solar Panels & - & 11.5W & 22W \\
GNSS Navigation & -0.21W & -0.255W & -0.39W \\
\midrule\noalign{}
\endhead
\bottomrule\noalign{}
\endlastfoot
\end{longtable}

*Payload data is yet to be obtained, Worse case scenarios for each state
is taken

\subsubsection{Communications System (COMMS)}

The Communications System provides the spacecraft's interface with the
ground segment and supports telemetry transmission etc. The COTS model
we have taken for the simulation is the GomSpace NanoCom AX100 radio
system working on two antenna systems: GomSpace ANT430 omnidirectional
antenna and ISIS UHF dipole antenna.

The primary ground station is planned for \textbf{Goa, India}.

\begin{longtable}{@{}
  >{\raggedright\arraybackslash}p{(\columnwidth - 2\tabcolsep) * \real{0.5957}}
  >{\raggedright\arraybackslash}p{(\columnwidth - 2\tabcolsep) * \real{0.4043}}@{}}
\caption{Goa ground station assumptions}\\
\toprule\noalign{}
\textbf{Parameter} & \textbf{Value} \\
Latitude & 15°N \\
Longitude & 74°E \\
Altitude & 70 m \\
Elevation mask & 15° \\
\midrule\noalign{}
\endhead
\bottomrule\noalign{}
\endlastfoot
\end{longtable}

Simulation predicts an average of \textbf{two contacts per day}, with
durations of approximately \textbf{368 s} and \textbf{414 s}. The
maximum elevations are approximately \textbf{38.5°} and \textbf{77.5°},
providing roughly \textbf{780 s of total daily contact time}.

Although our simulations currently have one ground station, we aim to
have more ground stations for receiving purposes.

\textbf{Communications Concept}

ANTARA uses a \textbf{Relay Mode} architecture, with communications
restricted to ground-station contact periods. During Relay Mode, the
spacecraft performs:

\begin{itemize}
\item
  Telemetry transmission
\item
  Science-data downlink
\item
  Command reception
\item
  Operational synchronisation
\end{itemize}

\textbf{Downlink Architecture}

The baseline downlink rate is \textbf{15--18 kbps}.

The scientific payload does not transmit individual proton events.
Instead, the onboard FPGA accumulates events into \textbf{64-bin
spatial-energy histograms}. Each histogram contains:

64×4=256 bytes

At one update per minute, the science data rate is approximately:

(256×8)/60≈34 bps

This corresponds to an estimated \textbf{300--400 KB/day} of science
data.

With two daily contacts totalling approximately \textbf{780 s}, a 15
kbps link provides:

15,000×780≈11.7 Mbit

Thus, the available capacity provides approximately a 4× margin over the
expected science-data volume, allowing for housekeeping, operational
logs, attitude data, protocol overhead, and retransmissions.

The \textbf{9.6 kbps} value used in earlier analysis is superseded by
the architecture-driven \textbf{15--18 kbps baseline}.

\subsubsection{Historical Payload Subsystem}
\noindent\textbf{Superseded architecture.} The silicon entrance, five-layer stopping telescope, rear veto, EJ-200 selection and calibrated-bin language below are historical only. Current construction and response limitations are in Appendix~\ref{app:payload-detail}.

\textbf{Overview}

The payload is a \textbf{Flux Particle Detector Unit (FPDU)}: a compact
proton spectrometer using a silicon ΔE front detector followed by a
5-layer EJ-200 plastic scintillator stack, read out by SiPMs, with a
rear veto layer. It measures trapped protons from \textbf{5--100 MeV} in
\textbf{64 log-spaced energy bins} at Sun-synchronous orbit
(\textasciitilde600 km)

\textbf{Physical Stack --- Dimensions Along Boresight}

\begin{longtable}{@{}
  >{\raggedright\arraybackslash}p{(\columnwidth - 4\tabcolsep) * \real{0.6147}}
  >{\raggedright\arraybackslash}p{(\columnwidth - 4\tabcolsep) * \real{0.1341}}
  >{\raggedright\arraybackslash}p{(\columnwidth - 4\tabcolsep) * \real{0.2512}}@{}}
\caption{FPDU telescope physical-stack dimensions}\\
\toprule\noalign{}
\textbf{Component} & \textbf{Thickness} & \textbf{Notes} \\
Aperture tube + structure & 30 mm & Sets θ\_max = 18.4° \\
Si wafer & 1 mm & Active detector \\
Si mount + PCB & 3 mm & Carrier board \\
 &  &  \\
L1 + SiPM + coupling & 4 mm & Scintillator** + SiPM package \\
L2 + SiPM + coupling & 6 mm &  \\
L3 + SiPM + coupling & 14 mm &  \\
L4 + SiPM + coupling & 24 mm &  \\
L5 + SiPM + coupling & 32 mm &  \\
Rear veto + SiPM & 4 mm &  \\
Optical isolation gaps & \textasciitilde3 mm & Black wrap between layers \\
Readout electronics board & 15 mm & FPGA, ADCs, shaping, bias \\
Structural margin & 5 mm &  \\
\textbf{Total} & \textbf{\textasciitilde141 mm} & Fits in 2U payload volume \\
\midrule\noalign{}
\endhead
\bottomrule\noalign{}
\endlastfoot
\end{longtable}

\textbf{**All these calculations are done assuming EJ200 plastic
scintillator and is subject to change.}

\textbf{Dead Time and Livetime}

\begin{longtable}{@{}
  >{\raggedright\arraybackslash}p{(\columnwidth - 2\tabcolsep) * \real{0.5398}}
  >{\raggedright\arraybackslash}p{(\columnwidth - 2\tabcolsep) * \real{0.4602}}@{}}
\caption{FPDU dead-time and livetime assumptions}\\
\toprule\noalign{}
\textbf{Parameter} & \textbf{Value} \\
\textbf{Dead time per event} & \textbf{\textasciitilde250 ns (shaping + ADC readout)} \\
\textbf{Peak total event rate (SAA core)} & \textbf{\textasciitilde100 Hz} \\
\textbf{Dead time fraction at peak flux} & \textbf{12.5\%} \\
\textbf{Livetime fraction at peak flux} & \textbf{87.5\%} \\
\textbf{Dead time at SAA periphery} & \textbf{\textless{} 0.1\% (negligible)} \\
\midrule\noalign{}
\endhead
\bottomrule\noalign{}
\endlastfoot
\end{longtable}

\subsubsection{Historical Attitude Determination and Control System (ADCS)}
\noindent\textbf{Superseded selection.} The CMG description below is not the current reaction-wheel and magnetorquer concept.

Attitude Determination and Control System(ADCS) continuously processes
sensor data to calculate the spacecraft\textquotesingle s exact position
relative to external reference points. The COTS model we have taken for
the simulation is TensorTech ADCS - 10m - D. The model uses a
comprehensive suite of components: one ADCS-Main Control Board
(ADCS-MCB), four TensorFSS-15M Fine Sun Sensors, and one integrated
Control Moment Gyroscope (CMG) model.

\begin{longtable}{@{}
  >{\raggedright\arraybackslash}p{(\columnwidth - 2\tabcolsep) * \real{0.4566}}
  >{\raggedright\arraybackslash}p{(\columnwidth - 2\tabcolsep) * \real{0.5434}}@{}}
\caption{TensorTech ADCS-10m-D preliminary characteristics}\\
\toprule\noalign{}
\textbf{Characteristic} & \textbf{Value} \\
Pointing Knowledge & 0.3 degree (general)

3-100 arcseconds (with startracker) \\
Pointing Accuracy & 0.4 degree \\
Power consumption & 0.8W (idle)

2.6W (nominal)

4.5W (Max) \\
Mass & 450g \\
\midrule\noalign{}
\endhead
\bottomrule\noalign{}
\endlastfoot
\end{longtable}

\textbf{Control Modes and Algorithms}

\textbf{Detumbling.} Following deployment, the ADCS executes a two-stage
control law: an initial B-dot magnetic detumble to reduce high angular
rates acquired during deployment, followed by reaction-wheel damping to
bring the spacecraft to a stable, low-rate condition suitable for Sun
Pointing Mode entry.

Simulation results show detumble completing in approximately 2,460
seconds from initial body rates of {[}6.0, −5.5, 4.5{]} °/s. No hard
detumble timeout is currently defined.

\textbf{Sun Pointing.} The ADCS orients the spacecraft +Z face toward
the Sun to maximize solar-array power generation, using reaction wheels
for fine pointing and magnetorquers for momentum dumping.

\textbf{Detection Mode Pointing.} During SAA traversal, the ADCS aligns
the +X axis (payload boresight) with the spacecraft velocity vector to
support directional proton flux measurement and East-West asymmetry
analysis. Simulation results show a mean attitude error of approximately
2°, against a threshold of ≤5° and a goal of ≤2°.

\textbf{Relay Mode Pointing.} The ADCS points the +Y face (directional
relay antenna) toward the ground station to optimize downlink signal
strength during scheduled contacts.

\textbf{Slewing.} Mode transitions between Sun Pointing, Detection, and
Relay require dedicated slew maneuvers. The 24-hour DITL simulation
recorded 18 slew events with slew angles ranging from 5.4° to 175.8°;
large-angle slews were identified as consuming non-negligible time and
should be factored into mission scheduling and duty-cycle analysis.

\textbf{Saturation Handling.} Reaction wheel momentum saturation is
managed through magnetorquer-based momentum unloading. If saturation
persists beyond a defined hold timeout and pointing diverges, the FSM
commands a transition to Safe Mode as a fault-protection response, this
is the second-highest-priority transition in the mode state machine
after critical battery conditions

\begin{longtable}{@{}
  >{\raggedright\arraybackslash}p{(\columnwidth - 2\tabcolsep) * \real{0.4760}}
  >{\raggedright\arraybackslash}p{(\columnwidth - 2\tabcolsep) * \real{0.5240}}@{}}
\caption{ADCS simulation performance summary}\\
\toprule\noalign{}
\textbf{Metric} & \textbf{Requirement / Result} \\
Detumble completion & \textasciitilde2,460 s from {[}6.0, −5.5, 4.5{]} °/s \\
Detection Mode pointing error & \textless2° \\
Slew events (24-hr sim) & 18 \\
Slew angle range & 5.4°--175.8° \\
Momentum management & Magnetorquer unloading + Safe Mode fallback \\
\midrule\noalign{}
\endhead
\bottomrule\noalign{}
\endlastfoot
\end{longtable}

Pointing Budget summary:

Ground station contact pointing budget analysis is being mentioned. The
result is for a simulation taken to Propagate 1-day Keplerian orbit
(analytical, no perturbations) with nadir-pointing attitude, computes
off-boresight angle between +Z body axis (antenna boresight) and ground
station LOS vector during passes above 10° elevation mask. Outputs pass
statistics and off-boresight angle plots.


\subsubsection{On-Board Computer (OBC)}

The On-Board Computer is the central computational and
operational-control unit of the ANTARA spacecraft. The COTS model we
have taken for the simulation is GomSpace Nanomind A3200.

\textbf{Purpose}

The OBC serves as the spacecraft's primary operational-management system
and continuously monitors spacecraft status while coordinating all major
mission activities.

At 3.3V bus channel, The OBC can run on various clock speeds from 8 -
32MHz depending on the state of the CubeSat with an average power
consumption of 0.17W

\textbf{Functional Responsibilities}

The ANTARA OBC performs mission sequencing, state monitoring, command
execution, health monitoring, payload control, data-storage management,
and communications scheduling. Payload control is closely integrated
with orbit propagation, ADCS orientation, SAA traversal prediction, and
spacecraft power availability, enabling synchronised science operations
during South Atlantic Anomaly crossings.

\textbf{Flight Software Philosophy}

ANTARA uses a state-machine-based flight software architecture designed
to maximise reliability, deterministic behaviour, operational
predictability, fault containment, and verification simplicity. Each
operational mode corresponds to a defined spacecraft state. State
transitions occur only when predefined operational conditions are
satisfied, such as battery thresholds, spacecraft attitude stability,
payload readiness, and subsystem health. The state-machine architecture
simplifies mission simulation, operational verification, fault
diagnosis, and long-duration autonomous operations.


\subsection{Spacecraft Coordinate System}

A spacecraft-fixed body-frame coordinate system is used throughout
ANTARA mission operations. The body-frame definition is designed around
payload operational geometry, solar-panel placement, communications
requirements, and CubeSat deployment constraints.

\begin{longtable}{@{}
  >{\raggedright\arraybackslash}p{(\columnwidth - 2\tabcolsep) * \real{0.1538}}
  >{\raggedright\arraybackslash}p{(\columnwidth - 2\tabcolsep) * \real{0.8462}}@{}}
\caption{Spacecraft body-frame configuration}\\
\toprule\noalign{}
\textbf{Axis} & \textbf{Function} \\
+X & RAM face / payload boresight --- primary payload pointing direction
during Detection Mode \\
+Y & Directional relay antenna face (nominally nadir-pointing in default
hold; pointed to ground station during Relay Mode) \\
+Z & Fixed solar panel face (pointed toward Sun during Sun Pointing Mode) \\
-Y & Deployable solar panel packed face \\
-Z & Omnidirectional communications antenna face (Safe Mode / beacon
transmission) \\
-X & PPOD deployment face \\
\midrule\noalign{}
\endhead
\bottomrule\noalign{}
\endlastfoot
\end{longtable}


\subsection{Ground Segment \& Mission Operations}

\textbf{Ground Station Hardware}

Primary ground station: Goa, India (15°N, 74°E, 70 m altitude, 15°
elevation mask).

Hardware baseline:

\begin{itemize}
\item
  \textbf{Antenna:} UHF Yagi or crossed-Yagi on an az-el tracking mount,
  matched to the NanoCom AX100 downlink band
\item
  \textbf{RF front-end:} transceiver/SDR compatible with AX100
  modulation, LNA at the antenna feed to minimize cable-loss impact on
  receive sensitivity
\item
  \textbf{Backup/contingency station:} none currently defined, remains
  an open item.
\end{itemize}

Secondary Ground Station: We are yet to determine secondary ground
stations. An ideal ground station would be at a polar location to
accommodate for the SSO orbit. A temporary simulation was taken for
Svalbard (78°N, 15°E, 50 m altitude, 5° elevation mask).

\begin{quote}
\textbf{Ground Station Software}
\end{quote}

\begin{itemize}
\item
  \textbf{Orbit propagation \& tracking:} TLE-driven antenna pointing;
  TLEs refreshed at minimum daily from CelesTrak/Space-Track, more
  frequently in the days immediately following deployment when orbit
  uncertainty is highest
\item
  \textbf{Pass prediction:} AOS/LOS scheduling consistent with the
  two-contacts-per-day, \textasciitilde780 s total daily contact time
\item
  \textbf{Command uplink interface:} command construction, queuing, and
  uplink receipt confirmation
\end{itemize}

\begin{itemize}
\item
  Documented shift handover: spacecraft state, open anomalies, and
  next-pass command queue passed between operators
\end{itemize}


